% Generated draft: requires mathematical and source review.
\documentclass[11pt]{article}
\usepackage[margin=25mm]{geometry}
\usepackage{fontspec}
\setmainfont{Latin Modern Roman}
\usepackage{amsmath,amssymb,mathtools,unicode-math}
\setmathfont{Latin Modern Math}
\usepackage{enumitem,needspace,xurl,hyperref}
\setlength{\parindent}{0pt}
\setlength{\parskip}{5pt}
\setlength{\emergencystretch}{3em}
\newcommand{\N}{\mathbb{N}}
\newcommand{\Z}{\mathbb{Z}}
\newcommand{\Q}{\mathbb{Q}}
\newcommand{\R}{\mathbb{R}}
\newcommand{\C}{\mathbb{C}}
\newcommand{\F}{\mathbb{F}}
\newcommand{\A}{\mathbb{A}}
\newcommand{\PP}{\mathbb{P}}
\newcommand{\E}{\mathbb{E}}
\newcommand{\eps}{\varepsilon}
\newcommand{\dd}{\,\mathrm{d}}
\begin{document}
\section*{UnsolvedMath ID: 6200014}
\textbf{Source problem code:} catalogue entry
\subsection{Definitions and mathematical statement}
Let \(N\) be a closed topological \(n\)-manifold admitting a finite triangulation \(\Delta\). A simplicial complex is flag if every finite set of pairwise adjacent vertices spans a simplex. The no-square condition conventionally means that its edge graph has no induced cycle of length four: every four-cycle has a diagonal.

For orientation, the standard Davis--Vinberg construction uses the right-angled Coxeter group
\[
 W_\Delta=\langle s_v\ (v\in\Delta^{(0)}):
 s_v^2=1,\quad s_vs_w=s_ws_v\text{ when }\{v,w\}\in\Delta^{(1)}\rangle
\]
and its Davis complex \(C(N,\Delta)\), with cubes corresponding to cosets of subgroups generated by cliques. Its visual boundary \(\partial_\infty C(N,\Delta)\) is the space of geodesic rays modulo bounded distance, with the visual topology for the intended proper nonpositively curved realization. The excerpt names the associated complex without specifying its precise metric model; agreement with this standard construction requires source confirmation.

Kapovich's question asks whether the boundary homeomorphism type depends only on the homeomorphism type of \(N\). In particular, for every closed \(N\) and every pair of flag, no-square triangulations \(\Delta_1,\Delta_2\) of \(N\), must
\[
 \partial_\infty C(N,\Delta_1)\cong
 \partial_\infty C(N,\Delta_2)
\]
hold? No equivariant or canonical homeomorphism is required by this formulation, only independence of the topological boundary type. No assertion that every manifold admits such a triangulation is added.

The supplied summary is internally unsettled: it labels the question open/triage and ``likely false'', but also categorically says boundaries are not independent in general. It refers to Januszkiewicz--Świątkowski and Davis--Januszkiewicz constructions without giving a counterexample meeting this exact closed-manifold, flag/no-square setup or a precise invariance theorem for the mentioned special classes. General dependence for other Davis complexes or hyperbolizations does not by itself identify a resolution of this question. These claims are therefore retained as unverified attributions rather than an answer.
\subsection{Short English statement}
Does the visual boundary of the Davis--Vinberg complex depend only on the closed manifold, rather than its flag, no-square triangulation? The summary mixes an open label with an unsupported general noninvariance claim, so the exact construction and any counterexample need source review.
\subsection{Sources}
\begin{itemize}
\item Local registry: \url{data/unsolvedmath/problems.json}, permanent UnsolvedMath ID \(6200014\), source code catalogue entry; supplied statement, background, and research summary, updated 2026-08-06.
\item Misha Kapovich, \emph{Problems on Boundaries of Groups and Kleinian Groups} (2005), Problem 14, PDF page 5: \url{https://www.math.ucdavis.edu/~kapovich/EPR/problems.pdf}. Supplied normalized PDF-text extraction accessed 2026-07-29; current status marked as needing review.
\item Januszkiewicz--Świątkowski and Davis--Januszkiewicz constructions are mentioned by the summary without publication details, exact hypotheses, or a specific counterexample.
\end{itemize}
\end{document}
